\documentclass[11pt]{exam}
\usepackage{foerderkurs}

\begin{document}

\werkzeug{W2}{Vertrauen in die eigene Mathematik}

«Ich kann einfach kein Mathe.» Diesen Satz hören Mathematiklehrpersonen oft.
Meistens stimmt er so nicht: Hinter dem Gefühl, Mathematik nicht zu können,
stecken fast immer konkrete Lücken, ungünstige Lerngewohnheiten oder schlechte
Erfahrungen. All das lässt sich ändern. Diese Einheit besteht aus drei Teilen,
die über das Jahr verteilt behandelt werden.

\begin{lehrperson}[Überblick]
\textbf{Aufbau:} drei Teile zu je ca.\ 15 Minuten, dazu Reflexionsaufträge, die
auch als Hausaufgabe möglich sind.

\begin{tabular}{@{}llp{7.2cm}@{}}
  Teil 1 & Lücke statt Talent        & zu Kursbeginn, sobald die Profile aus der Standortbestimmung vorliegen \\
  Teil 2 & Sich realistisch einschätzen & nach der ersten Fortschrittskontrolle (Okt./Nov.) \\
  Teil 3 & Fehler, Steckenbleiben, Prüfungen & vor der Prüfungszeit am Semesterende
\end{tabular}

\textbf{Grundhaltung:} Selbstvertrauen entsteht nicht durch Ermutigung,
sondern durch \emph{erlebten Erfolg}. Programme, die nur eine positive
Einstellung zum Lernen vermitteln («Growth Mindset»), zeigen in der Forschung
im Durchschnitt eher kleine Effekte. Diese Einheit gibt deshalb vor allem einen
\emph{Rahmen}, in dem die Schülerinnen und Schüler ihre Erfolge im Kurs
richtig deuten können: als Ergebnis von Übung und geschlossenen Lücken, nicht
als Zufall. Leere Beruhigungen («Alle können alles») vermeiden; ehrlich ist:
Nicht alle lernen gleich schnell, aber die Grundlagen können alle lernen.

\textbf{Unterrichtspraktiken, die die Einheit stützen:}
\begin{itemize}
  \item Rückmeldungen auf Strategie und Einsatz beziehen, nicht auf Talent:
        «Ihre Skizze hat den Weg gezeigt» statt «Sie sind ja begabt».
  \item Fortschritt sichtbar machen: Profile vergleichen, abgehakte Ziele.
  \item Anonyme Antwortformen (Mini-Whiteboards, Umfragen), damit auch
        unsichere Schülerinnen und Schüler antworten.
  \item Darauf achten, wer sich äussert, wer die anspruchsvollen Fragen
        erhält und wie lange auf Antworten gewartet wird.
\end{itemize}

\textbf{Junge Frauen:} Studien wie PISA zeigen, dass Schülerinnen ihre
Fähigkeiten in Mathematik bei gleicher Leistung im Durchschnitt tiefer
einschätzen. Die Einheit spricht das nicht als Geschlechterthema an, aber
Teil~2 macht es sichtbar: Wer viele Aufgaben «unsicher, aber richtig» löst,
unterschätzt sich. Das lässt sich im Einzelgespräch gut aufgreifen.

\textbf{Grenze:} Die Einheit ersetzt keine Beratung. Bei starker
Prüfungsangst oder wenn sich eine Schülerin oder ein Schüler sehr belastet
zeigt, auf die schulischen Beratungsangebote hinweisen.
\end{lehrperson}

% ===========================================================================
\werkzeugteil{1}{Lücke statt Talent}
  {erklären, warum sich Schwierigkeiten in Mathematik oft wie fehlendes
   Talent anfühlen, und in meinem Profil erkennen, welche Lücken ich
   schliessen kann.}

\begin{einstieg}
Wie sehr stimmen Sie zu? Kreuzen Sie an, ohne Ihren Namen aufzuschreiben.
\begin{center}
\renewcommand{\arraystretch}{1.5}
\begin{tabular}{p{9.5cm}cccc}
  & \small gar nicht & \small eher nicht & \small eher & \small völlig \\
  Manche Menschen sind einfach keine Mathe-Menschen.
    & $\Box$ & $\Box$ & $\Box$ & $\Box$ \\
  Wer etwas nicht sofort versteht, ist nicht begabt.
    & $\Box$ & $\Box$ & $\Box$ & $\Box$ \\
  Gute Mathematikerinnen und Mathematiker machen kaum Fehler.
    & $\Box$ & $\Box$ & $\Box$ & $\Box$ \\
  Wer schnell rechnet, ist gut in Mathematik.
    & $\Box$ & $\Box$ & $\Box$ & $\Box$
\end{tabular}
\end{center}
\end{einstieg}

\begin{lehrperson}[Teil 1, Einstieg]
Zettel einsammeln und die Verteilung kurz zeigen (Strichliste an der
Wandtafel). Meist stimmt ein grosser Teil der ersten Aussage zu. Nicht
moralisieren, sondern fragen: «Woran merkt man, dass jemand kein
Mathe-Mensch ist?» Die Antworten (z.\,B.\ «versteht es nicht sofort», «macht
viele Fehler») führen direkt zum Merkkasten. Zur letzten Aussage: Auch
in der Mathematik als Wissenschaft ist Schnelligkeit kaum wichtig, Ausdauer
dagegen sehr.
\end{lehrperson}

\begin{merke}
\begin{itemize}
  \item \textbf{Mathematik baut aufeinander auf}, wie ein Gebäude: Gleichungen
        stehen auf Termen, Terme auf Brüchen, Brüche auf dem Verständnis der
        Division. Fehlt ein Stockwerk weiter unten, wackelt alles darüber.
  \item Das fühlt sich an wie «Ich kann kein Mathe». Meistens ist es aber eine
        \textbf{konkrete Lücke}, oft viel weiter unten als das aktuelle Thema.
  \item \textbf{Lücken lassen sich schliessen.} Dafür ist dieser Kurs da. Ihre
        Standortbestimmung zeigt, wo Ihre Lücken liegen.
  \item Nicht sofort verstehen, langsam rechnen und Fehler machen sind
        \textbf{normale Teile des Lernens}, keine Zeichen fehlender Begabung.
\end{itemize}
\end{merke}

\begin{beispiel}[1: Wo liegt das Problem wirklich?]
Noah hat Mühe mit Gleichungen wie $\frac{x}{2} + \frac{x}{3} = 10$. Er sagt:
«Gleichungen kann ich einfach nicht.» Seine Standortbestimmung zeigt aber:

\begin{center}
\renewcommand{\arraystretch}{1.3}
\begin{tabular}{lcccc}
  & A & B & C & D \\
  3 Brüche      & \checkmark & \texttimes & \texttimes & \texttimes \\
  6 Gleichungen & \checkmark & \checkmark & \texttimes & \texttimes
\end{tabular}
\end{center}

Einfache Gleichungen löst Noah sicher (6B). Er scheitert bei Gleichungen mit
Brüchen, weil ihm das Rechnen mit Brüchen fehlt (3B). Sein Problem ist nicht
«Gleichungen», sondern eine Lücke bei den Brüchen. Und daran kann er gezielt
arbeiten.
\end{beispiel}

\begin{beispiel}[2: Wie erkläre ich mir Erfolg und Misserfolg?]
Nach einer misslungenen Prüfung sind ganz verschiedene Erklärungen möglich:

\begin{center}
\renewcommand{\arraystretch}{1.4}
\begin{tabular}{p{6.5cm}p{6.5cm}}
  \textbf{Erklärung} & \textbf{Kann ich das ändern?} \\
  «Ich bin halt nicht begabt.»            & nein, so formuliert nicht \\
  «Die Prüfung war unfair.»               & kaum \\
  «Ich habe zu wenig geübt.»              & ja \\
  «Ich habe nur die Lösungen gelesen, statt selbst zu rechnen.» & ja \\
  «Mir fehlen die Brüche.»                & ja
\end{tabular}
\end{center}

Erklärungen, die man beeinflussen kann, helfen weiterzumachen. Erklärungen,
die man nicht beeinflussen kann, führen dazu, dass man aufgibt, auch wenn das
gar nicht nötig wäre.
\end{beispiel}

\begin{questions}

\question Denken Sie an eine Mathematikprüfung, die nicht gut gelaufen ist.
Welche Erklärung hatten Sie damals? Ist diese Erklärung veränderbar? Finden
Sie eine andere, ehrliche Erklärung, die Sie beeinflussen können.
\begin{solution}
  Individuelle Antworten. Wichtig ist, dass die zweite Erklärung konkret ist,
  z.\,B.\ «Ich habe die Aufgaben mit Brüchen nicht geübt» statt «Ich habe zu
  wenig gelernt».
\end{solution}
\zeilen{4cm}

\question Schauen Sie Ihr Profil aus der Standortbestimmung an.
\begin{parts}
  \part In welchem Strang liegt Ihre tiefste Bruchstelle?
  \part Welche anderen Themen bauen auf diesem Strang auf?
  \part Formulieren Sie ein konkretes Ziel bis zur nächsten
        Fortschrittskontrolle, z.\,B.\ «Ich kann zwei Brüche mit
        verschiedenen Nennern ohne Fehler addieren.»
\end{parts}
\begin{solution}
  Individuelle Antworten. Gute Ziele sind konkret und überprüfbar. «Besser in
  Mathe werden» ist kein überprüfbares Ziel.
\end{solution}
\zeilen{5cm}

\end{questions}

\begin{transfer}
Auch in Sprachen, in Musik oder im Sport baut Können aufeinander auf. Wer in
Französisch die Verbformen nicht sicher beherrscht, hat bei jedem neuen
Grammatikthema Mühe und hält sich dann vielleicht für unbegabt. Die Frage
«Welches Stockwerk fehlt mir?» hilft überall.
\end{transfer}

\begin{exitfrage}
Ergänzen Sie den Satz: «Wenn ich in Mathematik etwas nicht verstehe, liegt
das oft daran, dass \ldots»
\ifprintanswers\par\medskip\textbf{Erwartet:} ein Bezug auf eine Lücke, auf
fehlende Übung oder eine Lernstrategie, nicht auf Begabung.
\else\par\vspace{1.5cm}\fi
\end{exitfrage}

% ===========================================================================
\newpage
\werkzeugteil{2}{Sich realistisch einschätzen}
  {vor dem Kontrollieren einschätzen, wie sicher ich bin, und daraus erkennen,
   wo ich mich über- oder unterschätze.}

\begin{einstieg}
Lösen Sie die Aufgaben. Kreuzen Sie \emph{vor} dem Kontrollieren an, wie
sicher Sie sind (1 = rate nur, 4 = ganz sicher). Erst danach kontrollieren
Sie mit den Lösungen.

\begin{center}
\renewcommand{\arraystretch}{2.2}
\newcommand{\sicher}{$\Box$\,1\quad$\Box$\,2\quad$\Box$\,3\quad$\Box$\,4}
\begin{tabular}{|l|p{3cm}|c|c|}
  \hline
  \textbf{Aufgabe} & \textbf{Ihr Ergebnis} & \textbf{Wie sicher?} & \textbf{richtig?} \\ \hline
  a) \ $\frac{2}{3} + \frac{1}{4}$          & \ifprintanswers$\frac{11}{12}$\fi & \sicher & \\ \hline
  b) \ $-3^2$                                & \ifprintanswers$-9$\fi            & \sicher & \\ \hline
  c) \ $20 \geteilt 4 \cdot 5$               & \ifprintanswers$25$\fi            & \sicher & \\ \hline
  d) \ $(x + 3)^2$                           & \ifprintanswers$x^2 + 6x + 9$\fi  & \sicher & \\ \hline
  e) \ $2x + 5 = 17$, \ $x = \;?$            & \ifprintanswers$6$\fi             & \sicher & \\ \hline
  f) \ $15\,\%$ von $60$                     & \ifprintanswers$9$\fi             & \sicher & \\ \hline
\end{tabular}
\end{center}
\end{einstieg}

\begin{lehrperson}[Teil 2, Einstieg]
Die Lösungen erst austeilen oder zeigen, wenn alle ihre Sicherheit
angekreuzt haben. Die Aufgaben sind so gewählt, dass typische
Fehlvorstellungen vorkommen (b, c, d): Dort entstehen oft «sicher, aber
falsch»-Fälle, die zu den wichtigsten Erkenntnissen dieser Lektion führen.
\end{lehrperson}

\begin{merke}
Sortieren Sie Ihre Aufgaben in die vier Felder:

\begin{center}
\renewcommand{\arraystretch}{1.5}
\begin{tabular}{|p{2.2cm}|p{5.4cm}|p{5.4cm}|}
  \hline
  & \textbf{richtig} & \textbf{falsch} \\ \hline
  \textbf{sicher} &
    \textbf{Gesichertes Wissen.} Gut so. &
    \textbf{Achtung!} Hier sitzt vermutlich eine falsche Regel fest. Diese
    Fehler sind am gefährlichsten, weil man sie nicht bemerkt. \\ \hline
  \textbf{unsicher} &
    \textbf{Sie können mehr, als Sie denken.} Hier fehlt Vertrauen, nicht
    Wissen. &
    \textbf{Bekannte Lücke.} Sie wissen, dass Sie hier üben müssen. Das ist
    ein guter Ausgangspunkt. \\ \hline
\end{tabular}
\end{center}

Wer sich \textbf{realistisch} einschätzt, lernt gezielter: Er übt dort, wo es
nötig ist, und verlässt sich dort, wo er es kann.
\end{merke}

\begin{beispiel}[Zwei Profile]
\textbf{Lara} kreuzt fast überall «4» an. Vier Aufgaben sind richtig, zwei
falsch (b und d). Lara überschätzt sich bei Potenzen und binomischen Formeln.
Sie würde diese Themen vor einer Prüfung nicht üben, weil sie sich sicher
fühlt.

\textbf{Elias} kreuzt fast überall «1» oder «2» an. Fünf Aufgaben sind richtig.
Elias unterschätzt sich. Er übt vor einer Prüfung alles gleich viel und ist
in der Prüfung unsicher, obwohl er vieles kann.

Beide profitieren davon, ihre Einschätzung mit dem Ergebnis zu vergleichen.
\end{beispiel}

\begin{questions}

\question Tragen Sie Ihre sechs Aufgaben aus dem Einstieg in die vier Felder
ein. In welchem Feld liegen die meisten? Was bedeutet das für Ihr Üben?
\begin{solution}
  Individuelle Antworten. Bei vielen «sicher, aber falsch»: diese Themen
  gezielt üben und die Regel überprüfen. Bei vielen «unsicher, aber richtig»:
  den eigenen Lösungen mehr vertrauen und das in Prüfungen bewusst nutzen.
\end{solution}
\zeilen{3cm}

\question Vergleichen Sie Ihr Profil aus der Standortbestimmung mit dem
Ergebnis der ersten Fortschrittskontrolle. Wo hat sich Ihre Bruchstelle
verschoben? Was hat dazu beigetragen?
\begin{solution}
  Individuelle Antworten. Wichtig ist der zweite Teil: Fortschritt soll auf
  konkretes Handeln zurückgeführt werden (bestimmte Übungen, eine
  Mini-Lektion, regelmässiges Aufwärmen).
\end{solution}
\zeilen{4cm}

\question Ab heute gilt bei jeder Übungsserie: Bevor Sie kontrollieren,
notieren Sie neben jeder Aufgabe eine Zahl von 1 bis 4. Führen Sie dazu ein
kurzes \textbf{Lernjournal}: Schreiben Sie am Ende jeder Kurslektion zwei
Sätze auf: Was habe ich heute gelernt? Wo bin ich noch unsicher?
\begin{solution}
  Kein Lösungsteil. Die Lehrperson sollte das Journal in den folgenden Wochen
  gelegentlich anschauen oder kurz besprechen, sonst schläft es ein.
\end{solution}

\end{questions}

\begin{transfer}
Die Frage «Wie sicher bin ich?» hilft in jedem Fach bei der
Prüfungsvorbereitung: Vokabeln, die man sicher zu können glaubt, aber nicht
kann, sind die gefährlichsten. Wer sich vor dem Kontrollieren festlegt,
merkt das.
\end{transfer}

\begin{exitfrage}
Welches der vier Felder ist am gefährlichsten, und warum?
\ifprintanswers\par\medskip\textbf{Lösung:} «sicher, aber falsch», weil man
den Fehler nicht bemerkt und darum auch nicht übt.
\else\par\vspace{1.5cm}\fi
\end{exitfrage}

% ===========================================================================
\newpage
\werkzeugteil{3}{Fehler, Steckenbleiben und Prüfungen}
  {meine Fehler einordnen, mir selbst weiterhelfen, wenn ich nicht
   weiterkomme, und mich wirksam auf eine Prüfung vorbereiten.}

\begin{einstieg}
Vier Schülerinnen und Schüler haben einen Fehler gemacht. Ist es jedes Mal
dieselbe Art von Fehler?
\begin{enumerate}
  \item Lena rechnet $7 \cdot 8 = 54$. Sie kann das Einmaleins eigentlich gut.
  \item Gefragt ist der Umfang eines Rechtecks mit den Seiten $6$\,cm und
        $4$\,cm. Jonas antwortet: $24$\,cm$^2$.
  \item Mira rechnet $\frac{2}{3} + \frac{1}{4} = \frac{3}{7}$ und ist sich
        ganz sicher.
  \item Bei der Aufgabe $(x + 3)^2$ schreibt Tim: «Ich weiss nicht, wie man
        das rechnet.»
\end{enumerate}
\end{einstieg}

\begin{lehrperson}[Teil 3, Einstieg]
Zu zweit 2 Minuten diskutieren lassen. Die Pointe: Alle vier Antworten sind
«falsch», aber die Ursachen sind völlig verschieden, und darum brauchen sie
verschiedene Gegenmittel. Wer nach einer Prüfung nur «Ich habe viele Fehler
gemacht» sagt, weiss nicht, was er ändern soll.
\end{lehrperson}

\begin{merke}
\textbf{Vier Arten von Fehlern}, und was jeweils hilft:

\begin{center}
\renewcommand{\arraystretch}{1.4}
\begin{tabular}{|p{3.3cm}|p{4.5cm}|p{5.3cm}|}
  \hline
  \textbf{Fehlerart} & \textbf{Merkmal} & \textbf{Was hilft} \\ \hline
  Flüchtigkeitsfehler & eigentlich gewusst & Kontrollroutine: Probe, Überschlag,
    am Schluss alles nochmals lesen \\ \hline
  Lesefehler & Frage falsch verstanden & Routine aus W1: gegeben und gesucht
    markieren \\ \hline
  Lücke & nicht gewusst, wie & gezielt üben, Mini-Lektion, nachfragen \\ \hline
  Fehlvorstellung & falsche Regel, sicher angewendet & Regel überprüfen und
    neu lernen; erkennbar an «sicher, aber falsch» \\ \hline
\end{tabular}
\end{center}
\end{merke}

\begin{merke}
\textbf{Wenn ich nicht weiterkomme:}
\begin{enumerate}
  \item Die Aufgabe nochmals lesen. Was ist gegeben, was ist gesucht?
  \item Eine Skizze oder eine Tabelle machen.
  \item Ein einfacheres Beispiel mit kleinen Zahlen ausprobieren.
  \item Eine ähnliche Aufgabe oder einen passenden Merkkasten suchen.
  \item Fragen: Was kann ich schon ausrechnen, auch wenn es nicht das Ziel ist?
  \item Aufschreiben, \emph{wo genau} ich hänge, und dann mit dieser genauen
        Frage um Hilfe bitten.
\end{enumerate}
Steckenbleiben ist kein Zeichen von Versagen. Es ist der Moment, in dem
Lernen stattfindet.
\end{merke}

\begin{merke}
\textbf{Prüfungen}

\emph{Vorbereitung:} Selbst Aufgaben lösen statt Lösungen lesen; über mehrere
Tage verteilt üben statt am Vorabend; eine Probeprüfung unter Zeitdruck lösen.

\emph{In der Prüfung:} Zuerst alle Aufgaben überfliegen; mit einer machbaren
Aufgabe beginnen; nicht an einer Aufgabe hängen bleiben, sondern später
zurückkehren; am Schluss Zeit für die Kontrolle einplanen.

\emph{Bei Nervosität:} Herzklopfen und Anspannung bedeuten, dass der Körper
sich auf eine Leistung vorbereitet. Langsam ausatmen hilft, ruhiger zu
werden. Manchen hilft es, vor der Prüfung ein paar Minuten lang ihre Sorgen
aufzuschreiben. Probieren Sie aus, was für Sie funktioniert.
\end{merke}

\begin{lehrperson}[Prüfungsangst]
Die Hinweise zur Nervosität stützen sich auf Studien zum Umdeuten von
Anspannung und zum Aufschreiben von Sorgen vor Prüfungen. Die Effekte sind
nicht bei allen gleich, darum die Formulierung «Probieren Sie aus». Wenn eine
Schülerin oder ein Schüler von starker, anhaltender Angst berichtet, das
Gespräch suchen und auf die Beratungsangebote der Schule hinweisen.
\end{lehrperson}

\begin{questions}

\question Lösen Sie den Einstieg: Ordnen Sie die vier Fehler einer Fehlerart
zu.
\begin{solution}
  1.\ Flüchtigkeitsfehler \quad 2.\ Lesefehler \quad 3.\ Fehlvorstellung
  \quad 4.\ Lücke
\end{solution}
\zeilen{2cm}

\question Nehmen Sie Ihre letzte Mathematikprüfung oder eine Übungsserie und
erstellen Sie ein Fehlerprotokoll.

\begin{center}
\renewcommand{\arraystretch}{2}
\begin{tabular}{|p{1.6cm}|p{4.2cm}|p{3.2cm}|p{4.2cm}|}
  \hline
  \textbf{Aufgabe} & \textbf{Was war falsch?} & \textbf{Fehlerart} &
  \textbf{Was tue ich dagegen?} \\ \hline
  & & & \\ \hline
  & & & \\ \hline
  & & & \\ \hline
  & & & \\ \hline
\end{tabular}
\end{center}
\begin{solution}
  Individuelle Antworten. Wiederkehrende Fehlerarten sind besonders
  aufschlussreich: Viele Flüchtigkeitsfehler deuten auf eine fehlende
  Kontrollroutine hin, viele Fehlvorstellungen auf einen bestimmten Strang,
  der in der Standortbestimmung ebenfalls auffällt.
\end{solution}

\question Wählen Sie aus der Liste «Wenn ich nicht weiterkomme» die drei
Schritte aus, die für Sie am hilfreichsten sind, und schreiben Sie sie auf
eine Karte, die Sie in Ihr Heft legen.
\begin{solution}
  Individuelle Antworten.
\end{solution}

\end{questions}

\begin{transfer}
Die vier Fehlerarten gibt es in jedem Fach: ein vergessenes Akzentzeichen in
Französisch (Flüchtigkeit), eine falsch verstandene Aufgabenstellung in
Geschichte (Lesefehler), eine nie gelernte Formel in Physik (Lücke), eine
falsche Vorstellung in Biologie (Fehlvorstellung). Und die
Prüfungsvorbereitung funktioniert überall gleich: selbst abrufen, verteilt
üben, unter Prüfungsbedingungen testen.
\end{transfer}

\begin{exitfrage}
Welche Fehlerart kommt bei Ihnen am häufigsten vor, und was werden Sie vor
der nächsten Prüfung dagegen tun?
\ifprintanswers\par\medskip\textbf{Erwartet:} eine konkrete Fehlerart und ein
konkretes Gegenmittel aus dem Merkkasten.
\else\par\vspace{1.5cm}\fi
\end{exitfrage}

\end{document}
